
Data curation is standard in LLM training pipelines, yet practitioners select filtering thresholds heuristically. This paper presents a controlled dose-response study of perplexity filtering thresholds for language model pretraining. Experiments on GPT-2 variants (125M to 1B parameters) using RedPajama-v2 data demonstrate that perplexity filtering exhibits a non-monotonic relationship with benchmark performance: intermediate thresholds near the 50th percentile outperform both extremes. A cubic polynomial fit achieves $R^2$ = 0.985 with an estimated optimal threshold at p53.9 (95\% CI: p40--p50). Convergence dynamics analysis shows unfiltered training exhibits 40\% higher convergence area-under-curve compared to moderate filtering, supporting the proposed noise-dilution mechanism. Scale transfer experiments indicate that optimal thresholds identified at 125M scale transfer to 1B scale with ratio 0.85. CPDR-optimized configuration shows 1.32\% improvement over RedPajama defaults at proof-of-concept scale. All experiments were conducted at reduced scale (5M--10M tokens) with synthetic or mock benchmark evaluation; effect magnitudes require full-scale replication for production deployment.
